\section*{Chapitre \ref{ch:iv1:timed-test-6} --- Épreuve chronométrée 6}
\markboth{\protect\omTocMark{}{Réponses~: Épreuve chronométrée 6}}{\protect\omTocMark{}{Réponses~: Épreuve chronométrée 6}}

\begin{solution}{iq:iv1:timed-test-6:1}
La tour attaque avec la probabilité $\frac29$, et le roi avec la probabilité $\frac{5}{48}$.

Une tour attaque les mêmes 14 cases où qu'elle se trouve~: 7 dans sa ligne et 7 dans sa colonne. L'autre pièce se trouve sur l'une des 63 cases restantes, donc la probabilité vaut $\frac{14}{63} = \frac29$.

La portée d'un roi dépend de l'endroit où il se trouve\index{conditionnement sur la position}. Il y a 4 coins avec 3 voisins, 24 cases de bord avec 5, et 36 cases intérieures avec 8. Les couples ordonnés de voisins totalisent $12 + 120 + 288 = 420$, sur $64 \times 63 = 4032$. La probabilité vaut donc $\frac{5}{48}$. Une énumération confirme les deux résultats.

\iqwrong{En prenant 8 voisins partout, ce qui donne $\frac{8}{63}$. On ignore que le bord de l'échiquier retranche environ un cinquième de ces voisins, puisque 28 des 64 cases en ont moins de 8.}
\end{solution}

\begin{solution}{iq:iv1:timed-test-6:2}
La première probabilité vaut $\frac12$, et la seconde $\frac38$.

Une équipe gagne ses 3 matchs avec la probabilité $\frac18$. Deux équipes ne peuvent pas le faire toutes deux, car elles se rencontrent. Ces événements sont donc disjoints\index{inclusion-exclusion}, et la somme vaut $4 \times \frac18$.

Pour la seconde, on choisit le vainqueur et le perdant, en $4 \times 3$ façons ordonnées. Ensemble, ces deux équipes disputent 5 matchs distincts, leur confrontation ne comptant qu'une fois. Chaque choix ordonné a donc la probabilité $2^{-5}$. Ces événements sont aussi disjoints, ce qui donne $\frac{12}{32} = \frac38$.

Une énumération sur les $2^6$ issues confirme les deux résultats. Avec $n$ équipes, la première probabilité vaut $n2^{-(n - 1)}$, qui tend vers 0 rapidement, comme il se doit.

\iqwrong{En prenant $1 - (\frac78)^4$, ce qui revient à traiter les quatre événements comme indépendants. Ils s'excluent mutuellement, de sorte que la somme simple est exacte.}
\end{solution}

\begin{solution}{iq:iv1:timed-test-6:3}
La position moyenne vaut $\frac{53}{5} = 10.6$, et la probabilité vaut $\frac{1}{17}$.

Les 4 as découpent les 48 autres cartes en 5 intervalles. Par symétrie, tous les intervalles ont la même taille moyenne\index{symétrie}, $\frac{48}{5} = 9.6$. Le premier as vient après le premier intervalle, si bien que sa position a pour moyenne $10.6$.

Quelle que soit la première carte, 3 des 51 cartes restantes ont le même rang qu'elle, et l'on a $\frac{3}{51} = \frac{1}{17}$.

Une simulation confirme les deux résultats. À la limite, avec 48 as et 4 autres cartes, la même règle donne $\frac{53}{49}$, juste au-dessus de 1, comme il se doit.

\iqwrong{En répondant 13, soit la taille du jeu divisée par le nombre d'as. C'est l'écart entre les as si le jeu est un cercle, et l'on ignore l'intervalle après le dernier as, qui prend lui aussi sa part.}
\end{solution}

\begin{solution}{iq:iv1:timed-test-6:4}
L'option d'achat vaut $\frac{15}{4} = 3.75$, répliquée par $\frac34$ d'action et $-33.75$ en liquidités. Une probabilité réelle de 90~\% ne change pas la valeur.

L'option d'achat rapporte 15 ou 0. En détenant $\Delta$ actions et $c$ en liquidités, il faut $65\Delta + c = 15$ et $45\Delta + c = 0$\index{réplication}. Donc $\Delta = \frac{15}{20} = \frac34$ et $c = -33.75$. Le coût vaut $\frac34 \times 50 - 33.75 = 3.75$.

Comme seconde méthode, la probabilité risque-neutre $p$ rend l'action équitable~: $65p + 45(1 - p) = 50$, donc $p = \frac14$, et la valeur vaut $\frac14 \times 15$.

Tout autre prix ouvre un arbitrage contre le portefeuille de réplication, quelles que soient les vraies cotes. Le gain espéré réel de $13.5$ n'est pas un prix.

\iqwrong{Valoriser au gain espéré réel. C'est l'erreur habituelle, et la réplication montre pourquoi elle échoue.}
\end{solution}

\begin{solution}{iq:iv1:timed-test-6:5}
Le nombre espéré de paires voisines de même couleur vaut 25, et le nombre espéré de paires dans une main vaut $\frac{10}{17}$.

Chacune des 51 positions voisines forme une paire de même couleur avec la probabilité $\frac{25}{51}$. Les variables indicatrices s'additionnent\index{variables indicatrices} sans qu'il soit besoin d'indépendance, si bien que la moyenne vaut exactement 25.

Dans une main, chacune des $\binom52 = 10$ paires de cartes a le même rang avec la probabilité $\frac{3}{51}$. La moyenne vaut donc $\frac{30}{51} = \frac{10}{17} \approx 0.59$.

Une simulation donne 25.0 et $0.588$. Notons qu'une main avec un brelan compte ici pour 3 paires, et que les variables indicatrices le comptent correctement.

\iqwrong{En prenant $\frac12$ pour la paire voisine, ce qui donne $25.5$. Une fois une carte rouge tirée, il reste une carte rouge de moins.}
\end{solution}

\begin{solution}{iq:iv1:timed-test-6:6}
Le nombre espéré de lignes vaut 2, et la probabilité vaut $\frac{37}{64}$.

Il y a 8 lignes. Chacune est monochrome avec la probabilité $2 \times 2^{-3} = \frac14$. Les variables indicatrices donnent\index{variables indicatrices} $8 \times \frac14 = 2$, bien que les lignes se recoupent.

Les trois rangées ne partagent aucune case, elles sont donc indépendantes. La probabilité qu'aucune rangée ne soit monochrome vaut $(\frac34)^3 = \frac{27}{64}$, et la réponse est le complémentaire.

Une énumération des 512 grilles confirme les deux résultats. Seules 32 des 512 grilles ne contiennent aucune ligne de ce type. La dépendance modifie la loi du nombre de lignes, mais pas sa moyenne.

\iqwrong{En prenant pour probabilité qu'il y ait au moins une ligne parmi les 8 la valeur $1 - (\frac34)^8$. Les lignes se partagent des cases, elles ne sont donc pas indépendantes. Seules les rangées le sont.}
\end{solution}

\begin{solution}{iq:iv1:timed-test-6:7}
Le nombre espéré vaut environ $3.49$, et la probabilité environ $0.0163$.

Une urne donnée est vide avec la probabilité $0.9^{10}$. La moyenne vaut donc $10 \times 0.349 = 3.49$\index{variables indicatrices}, un peu en dessous de la valeur limite pour $n$ grand, $10/e = 3.68$.

Avec exactement une urne vide, une urne contient deux boules et les huit autres en contiennent chacune une. On choisit l'urne vide en 10 façons. On choisit les deux boules qui partagent une urne, en $\binom{10}{2} = 45$ façons. On affecte les 9 groupes aux 9 urnes, en $9!$ façons. D'où
\[
P = \frac{10 \times 45 \times 362\,880}{10^{10}} = 0.0163.
\]

Une simulation confirme les deux valeurs. Avec $3.49$ urnes vides en moyenne, n'en avoir qu'une seule est rare, et c'est ce que traduit cette faible valeur.

\iqwrong{En multipliant $0.349$ par la probabilité que les autres urnes soient non vides, comme si les urnes étaient indépendantes. Remplir neuf urnes avec dix boules exige un dénombrement précis à la place.}
\end{solution}

\begin{solution}{iq:iv1:timed-test-6:8}
La probabilité vaut $\frac{244}{495} \approx 0.493$.

Le premier lancer gagne immédiatement avec la probabilité $\frac{6 + 2}{36} = \frac29$.

Pour un point obtenu de $w$ façons, seuls le point ou un 7 comptent aux lancers suivants. Le point sort donc en premier avec la probabilité $\frac{w}{w + 6}$\index{conditionnement sur l'événement pertinent}. Les points 4 et 10 ont $w = 3$, les points 5 et 9 ont $w = 4$, et les points 6 et 8 ont $w = 5$. Ainsi
\[
P = \frac29 + 2\left(\frac{3}{36} \cdot \frac39 + \frac{4}{36} \cdot \frac{4}{10} + \frac{5}{36} \cdot \frac{5}{11}\right).
\]
Le crochet vaut $\frac{1}{36} + \frac{2}{45} + \frac{25}{396}$. En additionnant le tout, on obtient $\frac{244}{495}$. Une simulation donne $0.493$. L'avantage de la maison est donc de $1.4~\%$ sur un pari à cote égale.

\iqwrong{En prenant $\frac{w}{36}$ pour la probabilité du point aux lancers suivants. Les lancers qui ne montrent ni le point ni un 7 ne comptent pas.}
\end{solution}

\begin{solution}{iq:iv1:timed-test-6:9}
La volatilité vaut environ $10.6~\%$, et la valeur limite est d'environ $8.9~\%$.

Pour $n$ actifs de même poids, de variance $\sigma^2$ et de corrélation $\rho$\index{diversification},
\[
\sigma_p^2 = \sigma^2\left(\frac{1}{n} + \frac{n - 1}{n}\rho\right).
\]
Avec $n = 10$, le crochet vaut $0.1 + 0.9 \times 0.2 = 0.28$. La variance vaut donc $0.04 \times 0.28 = 0.0112$, et la volatilité vaut $\sqrt{0.0112} = 10.6~\%$.

Quand $n$ croît, le crochet tend vers $\rho$, ce qui donne $\sqrt{0.008} = 8.9~\%$. C'est le risque que la diversification ne peut pas éliminer.

Une vérification aux limites~: $\rho = 1$ donne $20~\%$ quel que soit $n$, et $\rho = 0$ donne $20~\%/\sqrt{10} = 6.3~\%$.

\iqwrong{Diviser par $\sqrt{10}$ et s'arrêter à $6.3~\%$, ce qui suppose l'indépendance. La majeure partie du risque restant provient de la corrélation commune.}
\end{solution}

\begin{solution}{iq:iv1:timed-test-6:10}
La longueur moyenne vaut $\frac12 + x - x^2$, soit $\frac34$ en $x = \frac12$ et $\frac{11}{16}$ en $x = \frac14$.

Si la coupure $U$ tombe à gauche de $x$, le morceau est $(U, 1)$, de longueur $1 - U$. Si elle tombe à droite de $x$, le morceau est $(0, U)$, de longueur $U$. D'où
\[
\mathbb{E} = \int_0^x(1 - u)\,du + \int_x^1u\,du = x - \frac{x^2}{2} + \frac12 - \frac{x^2}{2}.
\]

Vérification aux limites~: en $x = 0$, le morceau commence toujours en 0, et sa longueur $U$ a pour moyenne $\frac12$, ce qui concorde. Une simulation donne les deux valeurs.

On remarque que le morceau qui contient un point fixe est plus long qu'un morceau moyen, de longueur $\frac12$. C'est l'échantillonnage biaisé par la longueur\index{echantillonnage@échantillonnage biaisé par la longueur}, qui sous-tend aussi le paradoxe de l'inspection.

\iqwrong{Répondre $\frac12$ pour tout $x$. Le point a plus de chances de tomber dans le plus long des deux morceaux.}
\end{solution}

\begin{solution}{iq:iv1:timed-test-6:11}
La longueur est uniforme sur $1, \dots, 10$, de moyenne $5.5$, et la probabilité vaut $\frac12$.

On suit l'élément 1 sous la permutation. Il revient sur lui-même dès le premier pas avec la probabilité $\frac{1}{10}$. Sinon, il le fait au pas suivant avec la probabilité $\frac19$, et ainsi de suite. La probabilité d'une longueur $k$ vaut
\[
\frac{9}{10} \cdot \frac89 \cdots \frac{11 - k}{12 - k} \cdot \frac{1}{11 - k} = \frac{1}{10}.
\]
La longueur est donc uniforme\index{structure cyclique}. Sachant que la longueur vaut $L$, l'élément 2 figure parmi les $L - 1$ autres membres du cycle avec la probabilité $\frac{L - 1}{9}$. Sa moyenne sur $L$ vaut $\frac{4.5}{9} = \frac12$.

Une simulation confirme ce résultat. Vérification aux limites~: pour $n = 2$, les deux éléments ne se retrouvent dans un même cycle que sous la transposition, de probabilité $\frac12$. La valeur ne dépend pas de $n$.

\iqwrong{Supposer que les longs cycles sont plus rares pour un élément donné. Une permutation compte plus de cycles courts que de longs, mais un long cycle contient beaucoup plus d'éléments, et les deux effets se compensent exactement.}
\end{solution}

\begin{solution}{iq:iv1:timed-test-6:12}
Lancer deux fois le dé, garder 21 des 25 issues, et recommencer sinon. Cela coûte $\frac{50}{21} \approx 2.38$ lancers par sortie, contre une borne de $\log_57 \approx 1.21$.

Deux lancers donnent 25 couples équiprobables. Attribuer à chacune des 7 sorties trois couples, soit 21 couples au total, et recommencer avec les 4 couples restants\index{méthode du rejet}. Chaque sortie a alors la même probabilité, si bien que le tirage est équilibré.

Chaque essai réussit avec la probabilité $\frac{21}{25}$, si bien que le nombre d'essais suit une loi géométrique de moyenne $\frac{25}{21}$. Chaque essai coûte 2 lancers. Une simulation renvoie $2.38$.

Chaque sortie porte $\log_27$ bits, et chaque lancer $\log_25$ bits, si bien qu'aucun schéma ne peut utiliser, en moyenne, moins de $\log_57$ lancers par sortie.

\iqwrong{Additionner deux lancers et réduire modulo 7. Les sommes ne sont pas équiprobables, donc la sortie n'est pas équilibrée, quelle que soit son apparence.}

\iqfollow{Pouvez-vous réutiliser les 4 issues rejetées~?}{Oui. Elles portent de l'information~: un couple rejeté est un tirage uniforme sur 4, qui peut alimenter la tentative suivante. Cela rapproche le coût de la borne.}
\end{solution}
